\documentclass[11pt]{article}
\usepackage[margin=1in]{geometry}
\usepackage{amsmath,amssymb,amsthm}
\usepackage{graphicx}
\usepackage{booktabs}
\usepackage{xcolor}
\usepackage{listings}
\usepackage[numbers]{natbib}
\usepackage{hyperref}

\newtheorem{proposal}{Proposed statement}
\theoremstyle{remark}
\newtheorem*{remark}{Remark}

\emergencystretch=2em

\lstdefinelanguage{lean}{
  morekeywords={def,theorem,example,namespace,end,open,fun,Prop,by,noncomputable,section,import},
  sensitive=true,
  morecomment=[l]{--},
  morecomment=[s]{/-}{-/},
}
\lstset{
  language=lean,
  basicstyle=\ttfamily\footnotesize,
  keywordstyle=\color{blue!60!black}\bfseries,
  commentstyle=\color{gray},
  columns=fullflexible,
  keepspaces=true,
  breaklines=true,
  frame=single,
  extendedchars=true,
  literate=
    {ℝ}{{$\mathbb{R}$}}1
    {ℕ}{{$\mathbb{N}$}}1
    {→}{{$\to$}}1
    {↔}{{$\leftrightarrow$}}1
    {∀}{{$\forall$}}1
    {∃}{{$\exists$}}1
    {∈}{{$\in$}}1
    {∑}{{$\sum$}}1
    {∏}{{$\prod$}}1
    {∧}{{$\wedge$}}1
    {≤}{{$\le$}}1
    {≠}{{$\ne$}}1
    {⨆}{{$\bigsqcup$}}1
    {ρ}{{$\rho$}}1
    {⟨}{{$\langle$}}1
    {⟩}{{$\rangle$}}1
}

\title{Proposed Sealed Statements for a Circular-Economy Kernel:\\
Loop Mass Balance, Cascaded Recovery, a Recovery-Adjusted Work Floor,\\
and the $N$-Partner Symbiosis Condition}
\author{Justin Hart, Viridis LLC}
\date{5 November 2026}

\begin{document}
\maketitle

\begin{abstract}
A software kernel for the circular economy has to do a few pieces of arithmetic correctly every
time. It must split demand into virgin and recovered input, track a secondary stock, compose
recovery stages, bound the fresh work needed when recovered supply is available, and decide
when a ring of firms exchanging by-products can leave every member strictly better off. We
propose four formal target statements for this arithmetic, written in Lean~4 over Mathlib, as
candidates for sealing in the Viridis Research program. \emph{LoopMassBalance} gives the
virgin/recovered split, an exact geometric closed form and the convergence of the secondary
stock, and strict monotonicity of cumulative virgin input in the recovery fraction.
\emph{CascadedRecovery} bounds the virgin share $\prod_i(1-r_i)$ after $n$ stages from below by
$(1-\max_i r_i)^n$ and from above by every single stage's share, strictly so under a stated
condition. \emph{RecoveryWorkFloor} recasts a canon work-ledger bound as a floor on fresh supply
given recovered supply and shows that the floor is attained. \emph{CyclicSymbiosis} states that a
strictly positive by-product flow around an $N$-partner cycle giving every partner strictly
positive net benefit exists if and only if $\prod_i c_i<\prod_i v_i$. Its $N=2$ case coincides
with the product condition of the two-partner Reciprocity Corridor. Every module compiles under
the pinned toolchain and carries concrete sanity checks and negative controls. Simple automated
tactics close none of the four statements. We make no claim to have proved any proposal. The
statements fix the mathematics of a model and say nothing about any real recovery rate, cost,
value or material.
\end{abstract}

\section{Introduction}
Circular-economy accounting rests on simple identities and inequalities. Material that is
recovered does not have to be extracted. Recovery stages in series compound. Recovered energy
or work offsets fresh supply. A ring of firms exchanging by-products is viable only if the
exchange creates value overall
\citep{geissdoerfer2017,chertow2000}. These statements are elementary, but a decision kernel that
applies them automatically should rely on exact, machine-checked versions with explicitly listed
hypotheses, not on informal folklore. Edge cases are where silent errors creep in. Examples are
full stock retention, out-of-range recovery fractions, stages that recover nothing, and partners
whose own exchange is unprofitable but whose cycle is profitable.

Viridis Research separates \emph{proposing} a formal target from \emph{proving} it. A proposal
is a Lean module that defines a proposition \texttt{Statement}. If the module compiles and simple
tactics do not close it, the program owner decides whether to seal it as an exact target for
later proof sessions. This paper is the first entry of project P3 (``Circular-economy kernel'').
It drafts four such modules and explains the design choices behind them. The question we address
is: \emph{what exact statements should the circular-economy kernel's formal core consist of?}
Each statement must be faithful to the intended model, strong enough to be useful, provably
non-vacuous, and continuous with the two canon results it builds on.

\paragraph{Contributions.}
\begin{itemize}
\item Four self-contained Lean~4 statement modules over Mathlib
\citep{moura2021lean4,mathlib2020} (Section~\ref{sec:statements}), each with sanity witnesses and
negative controls checked at compile time.
\item A recovery-adjusted reading of the canon work-ledger bound \citep{hart2026udwl}, with an
attainment half that the canon lacks.
\item An $N$-partner generalisation of the canon Reciprocity Corridor's feasibility theorem
\citep{hart2026reciprocity}, together with a compile-time check that its $N=2$ instance reproduces
the canon's \texttt{ProductCondition} and payoffs.
\item An automated probe (Section~\ref{sec:results}) showing that none of the four statements is
closed by \texttt{simp}, \texttt{norm\_num}, \texttt{decide}, \texttt{aesop}, \texttt{trivial},
\texttt{simp\_all} or \texttt{exact?}.
\end{itemize}

\section{Background and related work}
\paragraph{Circular economy and industrial symbiosis.}
\citet{geissdoerfer2017} review how the circular-economy concept relates to sustainability. Its
operational core is closing material loops, which means substituting recovered for virgin
material. \citet{chertow2000} defined industrial symbiosis as the exchange of materials, energy,
water and by-products among co-located firms and gave a taxonomy of exchange types. Our
\emph{CyclicSymbiosis} model is a deliberately stylised, linear version of one such exchange
pattern: a single directed ring of by-product transfers.

\paragraph{Canon results.}
The Viridis canon contains two results that P3 builds on. The Reciprocity Corridor (Run-128)
\citep{hart2026reciprocity} studies two partners $X$, $Y$ with payoffs $u_X = v_X y - c_X x$ and
$u_Y = v_Y x - c_Y y$. Its Lean formalisation proves that strictly positive flows with
$u_X,u_Y>0$ exist if and only if $c_Xc_Y<v_Xv_Y$ (\texttt{ProductCondition}). The Useful and
Dissipated Work Ledger (Run-163) \citep{hart2026udwl} defines useful work
$e_1w_1+e_2w_2$ from two supplies and proves \texttt{target\_requires\_supply}: if every
efficiency is at most $e_{\max}>0$ and useful work reaches $T$, then $w_1+w_2\ge T/e_{\max}$.

\paragraph{Formal tools.}
All modules are written for Lean~4 \citep{moura2021lean4} (toolchain v4.28.0) and Mathlib
\citep{mathlib2020} (commit \texttt{8f9d9cff}). They use only standard Mathlib notions: finite
products and sums, \texttt{Filter.Tendsto}, \texttt{StrictAntiOn} and \texttt{iSup} on~$\mathbb{R}$.

\section{Design principles}
\label{sec:design}
\begin{enumerate}
\item \textbf{Self-contained.} Each module imports only Mathlib, so a sealed statement does not
depend on canon files that might change. Where a canon definition is needed for comparison
(\texttt{ProductCondition}), it is copied verbatim.
\item \textbf{All hypotheses explicit.} Each hypothesis appears in the proposition as an
implication. Nothing is hidden in a structure or type class.
\item \textbf{Faithful to the brief, slightly stronger where cheap.} For example, LoopMassBalance
asserts the exact closed form for the gap $S(t)-S^\ast$ as well as convergence, and
CascadedRecovery adds a uniform-upper-bound version of the max bound.
\item \textbf{Controls compiled with the module.} Every module contains \texttt{example}s with
concrete numbers, proved by \texttt{norm\_num}, \texttt{simp} or \texttt{decide}. They include
\emph{witnesses}, which show that the hypotheses are satisfiable and the conclusions
non-trivial, and \emph{negative controls}, which show that dropping a hypothesis breaks a
conclusion.
\item \textbf{No field claims.} The statements are about models. Recovery rates, retention
rates, efficiencies, costs and values are free real parameters.
\end{enumerate}

\section{The four proposed statements}
\label{sec:statements}

\subsection{P3.LoopMassBalance}
\begin{proposal}[Loop mass balance]
Let $D\ge0$, $r\in[0,1]$, $\rho\in[0,1)$ and $S(0)\in\mathbb{R}$. Put $V=D(1-r)$,
$S(t+1)=\rho S(t)+rD$ and $S^\ast=rD/(1-\rho)$. Then $0\le V\le D$;
$S(t)-S^\ast=\rho^t(S(0)-S^\ast)$ for all $t\in\mathbb{N}$; $S(t)\to S^\ast$; and
$\sum_{t<T}V = TD(1-r)$. Moreover, if $D>0$ and $T>0$, then $r\mapsto\sum_{t<T}V$ is strictly
decreasing on $[0,1]$.
\end{proposal}
\emph{Informal argument.} The gap $g_t = S(t)-S^\ast$ satisfies $g_{t+1}=\rho g_t$ because
$\rho S^\ast + rD = S^\ast$. Since $0\le\rho<1$, $\rho^t\to0$. The remaining parts are
elementary. (No Lean proof is claimed here.)

\begin{lstlisting}[caption={Core of \texttt{proposals/P3.LoopMassBalance.lean} (docstrings omitted; \texttt{nhds} is the name behind Mathlib's $\mathcal{N}$ notation used in the file).}]
namespace ViridisStatements.P3.LoopMassBalance

open Filter Topology

def virgin (D r : ℝ) : ℝ := D * (1 - r)

def stock (ρ r D S0 : ℝ) : ℕ → ℝ
  | 0 => S0
  | t + 1 => ρ * stock ρ r D S0 t + r * D

def steadyStock (ρ r D : ℝ) : ℝ := r * D / (1 - ρ)

def cumulativeVirgin (D r : ℝ) (T : ℕ) : ℝ := ∑ _t ∈ Finset.range T, virgin D r

def Statement : Prop :=
  (∀ D r ρ S0 : ℝ, 0 ≤ D → 0 ≤ r → r ≤ 1 → 0 ≤ ρ → ρ < 1 →
    (0 ≤ virgin D r ∧ virgin D r ≤ D) ∧
    (∀ t : ℕ, stock ρ r D S0 t - steadyStock ρ r D
        = ρ ^ t * (S0 - steadyStock ρ r D)) ∧
    Tendsto (stock ρ r D S0) atTop (nhds (steadyStock ρ r D)) ∧
    (∀ T : ℕ, cumulativeVirgin D r T = (T : ℝ) * D * (1 - r))) ∧
  (∀ (D : ℝ) (T : ℕ), 0 < D → 0 < T →
    StrictAntiOn (fun r => cumulativeVirgin D r T) (Set.Icc 0 1))
\end{lstlisting}

\emph{Controls.} The witness $D=10$, $r=\rho=\tfrac12$, $S(0)=0$ gives $V=5$, the stock sequence
$0,5,\tfrac{15}{2},\dots$ and $S^\ast=10$, and the gap at $t=2$ equals
$(\tfrac12)^2(0-10)$. Over three periods, cumulative virgin input falls from $15$ to $\tfrac{15}{2}$
when $r$ rises to $\tfrac34$. There are three negative controls. With $\rho=1$ the stock grows
without bound and the closed form fails; Lean's convention $x/0=0$ makes $S^\ast=0$. With $D=0$
the cumulative virgin input does not depend on $r$. With $r=2$ the virgin input is negative.

\subsection{P3.CascadedRecovery}
\begin{proposal}[Cascaded recovery]
Let $n\in\mathbb{N}$ and $r_1,\dots,r_n\in[0,1]$, and let $\sigma=\prod_i(1-r_i)$ be the virgin
share. Then $0\le\sigma\le1$, and $(1-m)^n\le\sigma$ for every $m\le 1$ with $r_i\le m$ for all
$i$. In particular $(1-\sup_i r_i)^n\le\sigma$. Also $\sigma\le 1-r_j$ for every $j$, with strict
inequality whenever $r_j<1$ and $r_k>0$ for some $k\ne j$.
\end{proposal}
\emph{Informal argument.} Each factor lies in $[1-m,1]\subseteq[0,1]$. Splitting off factor $j$
gives $\sigma=(1-r_j)\prod_{i\ne j}(1-r_i)$, and the remaining product is at most $1-r_k<1$.
The maximum is written as Mathlib's $\bigsqcup_i r_i$ (\texttt{iSup}). On $\mathrm{Fin}\,0$ this
is $\sup\emptyset=0$, and both sides then equal~$1$, so no side condition $n>0$ is needed.

\begin{lstlisting}[caption={Core of \texttt{proposals/P3.CascadedRecovery.lean}.}]
namespace ViridisStatements.P3.CascadedRecovery

def virginShare {n : ℕ} (r : Fin n → ℝ) : ℝ := ∏ i, (1 - r i)

def Statement : Prop :=
  ∀ (n : ℕ) (r : Fin n → ℝ), (∀ i, 0 ≤ r i) → (∀ i, r i ≤ 1) →
    (0 ≤ virginShare r ∧ virginShare r ≤ 1) ∧
    (∀ m : ℝ, (∀ i, r i ≤ m) → m ≤ 1 → (1 - m) ^ n ≤ virginShare r) ∧
    (1 - ⨆ i, r i) ^ n ≤ virginShare r ∧
    (∀ j, virginShare r ≤ 1 - r j) ∧
    (∀ j k, k ≠ j → 0 < r k → r j < 1 → virginShare r < 1 - r j)
\end{lstlisting}

\emph{Controls.} With $r=(\tfrac12,\tfrac12)$ the share is $\tfrac14$, which meets the lower
bound and lies strictly below $\tfrac12$. With $r=(\tfrac12,\tfrac14,0)$ the share is
$\tfrac38>\tfrac18$. Four negative controls show which hypotheses matter. With $r=(0,\tfrac12)$
the cascade equals stage~2, so strictness needs another stage with $r_k>0$. With
$r=(\tfrac12,1)$ the share is $0=1-r_2$, so strictness needs $r_j<1$. With $r=(2,\tfrac12)$ the
share is negative. With $r=(-1,\tfrac12)$ the share exceeds $1-r_2$.

\subsection{P3.RecoveryWorkFloor}
\begin{proposal}[Recovery-adjusted work floor]
Let $w_f,w_r\ge0$, $e_f,e_r\in[0,e_{\max}]$ with $e_{\max}>0$, and $U=e_fw_f+e_rw_r$. If $U\ge T$,
then $w_f\ge (T-e_{\max}w_r)/e_{\max}$. Conversely, if $e_{\max}w_r\le T$, then
$w_f^\ast=(T-e_{\max}w_r)/e_{\max}$ is non-negative and, at $e_f=e_r=e_{\max}$, delivers exactly $T$.
\end{proposal}
\emph{Informal argument.} We have $T\le U\le e_{\max}(w_f+w_r)$. Dividing by $e_{\max}$ and
rearranging gives the floor. Attainment is a direct calculation.

\begin{remark}
Since $(T-e_{\max}w_r)/e_{\max}=T/e_{\max}-w_r$, the floor half is algebraically equivalent to
Run-163's \texttt{target\_requires\_supply} specialised to a fresh and a recovered supply. The
proposal is therefore a conservative extension. What it adds is the recovery reading (one unit
of recovered supply at the cap offsets one unit of fresh supply) and the attainment half, which
shows that the floor cannot be raised. Separate caps for fresh and recovered efficiencies would
be a genuinely stronger alternative, and we flag it for the sealing decision.
\end{remark}

\begin{lstlisting}[caption={Core of \texttt{proposals/P3.RecoveryWorkFloor.lean}.}]
namespace ViridisStatements.P3.RecoveryWorkFloor

def useful (wf wr ef er : ℝ) : ℝ := ef * wf + er * wr

def freshFloor (T wr emax : ℝ) : ℝ := (T - emax * wr) / emax

def Statement : Prop :=
  (∀ wf wr ef er emax T : ℝ, 0 ≤ wf → 0 ≤ wr →
    0 ≤ ef → ef ≤ emax → 0 ≤ er → er ≤ emax → 0 < emax →
    T ≤ useful wf wr ef er → freshFloor T wr emax ≤ wf) ∧
  (∀ wr emax T : ℝ, 0 < emax → 0 ≤ wr → emax * wr ≤ T →
    0 ≤ freshFloor T wr emax ∧ useful (freshFloor T wr emax) wr emax emax = T)
\end{lstlisting}

\emph{Controls.} The Run-163 witness $(8,4,\tfrac34,\tfrac12)$ gives $U=8$ and floor
$\tfrac{20}{3}\le8$. For attainment, $e_{\max}=\tfrac12$, $w_r=4$ and $T=8$ give floor $12$, and
$12$ units of fresh supply deliver exactly $8$. Without recovery the floor would be $16$. There
are three negative controls. Fresh supply $11<12$ falls short of the target. An efficiency of $1$
above the cap reaches the target below the floor. With $e_{\max}w_r>T$ the floor is negative.

\subsection{P3.CyclicSymbiosis}
\begin{proposal}[$N$-partner cyclic symbiosis]
Let $N\ge2$ and $c_i,v_i>0$ for $i\in\mathbb{Z}/N$. Partner $i$ ships $x_i$ at cost $c_ix_i$ and
receives value $v_ix_{i-1}$. A flow $x$ with all $x_i>0$ and all net benefits
$v_ix_{i-1}-c_ix_i>0$ exists if and only if $\prod_i c_i<\prod_i v_i$.
\end{proposal}
\emph{Informal argument.} ($\Rightarrow$) Multiply the $N$ strict inequalities
$c_ix_i<v_ix_{i-1}$ between positive numbers. Index shift is a bijection, so
$\prod_i x_{i-1}=\prod_i x_i>0$ cancels. ($\Leftarrow$) Put $a_i=v_i/c_i$, $P=\prod a_i>1$ and
$\lambda=P^{1/N}>1$. Set $x_0=1$ and $x_i=x_{i-1}a_i/\lambda$ for $1\le i<N$. Then
$x_i<a_ix_{i-1}$ for $i\ge1$, and $a_0x_{N-1}=P/\lambda^{N-1}=\lambda>1=x_0$ closes the cycle.

\begin{lstlisting}[caption={Core of \texttt{proposals/P3.CyclicSymbiosis.lean}. \texttt{ProductCondition} is copied verbatim from the canon file.}]
namespace ViridisStatements.P3.CyclicSymbiosis

def prev {N : ℕ} (i : Fin N) : Fin N :=
  ⟨(i.val + (N - 1)) % N, Nat.mod_lt _ (Nat.zero_lt_of_lt i.isLt)⟩

def net {N : ℕ} (c v x : Fin N → ℝ) (i : Fin N) : ℝ := v i * x (prev i) - c i * x i

def StrictlyBeneficialFlow {N : ℕ} (c v x : Fin N → ℝ) : Prop :=
  (∀ i, 0 < x i) ∧ ∀ i, 0 < net c v x i

def CycleProductCondition {N : ℕ} (c v : Fin N → ℝ) : Prop := ∏ i, c i < ∏ i, v i

def ProductCondition (vX cX vY cY : ℝ) : Prop := cX * cY < vX * vY

def Statement : Prop :=
  ∀ (N : ℕ), 2 ≤ N → ∀ c v : Fin N → ℝ, (∀ i, 0 < c i) → (∀ i, 0 < v i) →
    ((∃ x : Fin N → ℝ, StrictlyBeneficialFlow c v x) ↔ CycleProductCondition c v)
\end{lstlisting}

\paragraph{Agreement with the canon at $N=2$.}
The module checks two facts at compile time. First,
\texttt{CycleProductCondition ![cX, cY] ![vX, vY]} is equivalent to Run-128's
\texttt{ProductCondition vX cX vY cY}. Second, at the flow \texttt{![x, y]} partner~0 nets
$v_Xy-c_Xx=u_X$ and partner~1 nets $v_Yx-c_Yy=u_Y$. Hence the $N=2$ instance of the statement is
the canon's strict-feasibility theorem in its bilateral-benefit form.

\begin{lstlisting}[caption={The $N=2$ agreement checks, as they appear in the module.}]
example (vX cX vY cY : ℝ) :
    CycleProductCondition ![cX, cY] ![vX, vY] ↔ ProductCondition vX cX vY cY := by
  simp [CycleProductCondition, ProductCondition, Fin.prod_univ_two]

example (vX cX vY cY x y : ℝ) :
    net ![cX, cY] ![vX, vY] ![x, y] 0 = vX * y - cX * x ∧
    net ![cX, cY] ![vX, vY] ![x, y] 1 = vY * x - cY * y := by
  have h0 : prev (0 : Fin 2) = 1 := by decide
  have h1 : prev (1 : Fin 2) = 0 := by decide
  simp [net, h0, h1]
\end{lstlisting}

\emph{Controls.} The values of \texttt{prev} on $\mathrm{Fin}\,3$ and $\mathrm{Fin}\,2$ are
checked by \texttt{decide}. The three-partner economy $c=(1,1,1)$, $v=(4,\tfrac12,1)$ satisfies
$1<2$ even though partner~1 alone has $v_1<c_1$. The uniform flow gives partner~1 a negative
net, while $x=(1,\tfrac{9}{20},\tfrac25)$ gives nets $(\tfrac35,\tfrac1{20},\tfrac1{20})$, all
positive. The flow must therefore be tuned to the cycle. There are two negative controls.
$c=(2,2,2)$, $v=(1,1,1)$ fails the condition, and the uniform flow harms partner~0. The
boundary case $c=v=(1,1)$ fails the strict condition.

\section{Results}
\label{sec:results}
All four modules compile under the pinned toolchain without errors or warnings, including every
sanity example. We probed each \texttt{Statement} with seven simple tactics
(\texttt{code/probe.sh}, output in \texttt{results/probe\_results.tsv}). Table~\ref{tab:probe}
summarises the outcome. No tactic closed any statement, so none is trivial in the sense the
verifier tests. Table~\ref{tab:hyp} lists the hypotheses of each statement. These results describe
the proposals only. Proofs belong to later tasks, after a sealing decision.

\begin{table}[h]
\centering
\begin{tabular}{lccccccc}
\toprule
Statement & \texttt{simp} & \texttt{norm\_num} & \texttt{decide} & \texttt{aesop} & \texttt{trivial} & \texttt{simp\_all} & \texttt{exact?}\\
\midrule
LoopMassBalance & open & open & open & open & open & open & open \\
CascadedRecovery & open & open & open & open & open & open & open \\
RecoveryWorkFloor & open & open & open & open & open & open & open \\
CyclicSymbiosis & open & open & open & open & open & open & open \\
\bottomrule
\end{tabular}
\caption{Probe outcomes. ``open'' means the tactic failed to close the statement.}
\label{tab:probe}
\end{table}

\begin{table}[h]
\centering
\small
\begin{tabular}{lp{10.5cm}}
\toprule
Statement & Hypotheses \\
\midrule
LoopMassBalance & $0\le D$, $0\le r\le1$, $0\le\rho<1$ (parts 1--4); $0<D$, $0<T$, $r\in[0,1]$ (strict monotonicity)\\
CascadedRecovery & $0\le r_i\le1$ for all $i$; $r_i\le m\le1$ (uniform bound); $k\ne j$, $r_k>0$, $r_j<1$ (strictness)\\
RecoveryWorkFloor & $w_f,w_r\ge0$, $0\le e_f,e_r\le e_{\max}$, $e_{\max}>0$, $T\le U$ (floor); $e_{\max}>0$, $w_r\ge0$, $e_{\max}w_r\le T$ (attainment)\\
CyclicSymbiosis & $N\ge2$, $c_i>0$, $v_i>0$ for all $i$\\
\bottomrule
\end{tabular}
\caption{Hypotheses of each proposed statement.}
\label{tab:hyp}
\end{table}

\section{Limitations}
\begin{itemize}
\item \textbf{Models, not measurements.} No statement says anything about real recovery rates,
retention, efficiencies, costs, values, materials or environmental outcomes. Those remain field
inputs.
\item \textbf{Stylised dynamics.} LoopMassBalance has constant demand and rates, a single
material and no feedback from stock to recovery. CascadedRecovery treats stages as independent.
CyclicSymbiosis covers a single directed cycle with linear payoffs and no capacities, transport
losses or division of surplus. Run-128's Nash-midpoint result has no $N$-partner analogue here.
\item \textbf{Mild extension in RecoveryWorkFloor.} Its floor half is a rearrangement of the
canon bound (see the remark above). It is not a thermodynamic bound: it is neither a Landauer
bound nor dependent on the Intelligence Bound, which remains a conjecture.
\item \textbf{Probes are not proofs of hardness.} Failing seven simple tactics does not show
that a statement is deep. It only shows that the statement is not closed by those tactics.
\item \textbf{Tooling.} An attempt to use the Aristotle prover failed at submission with an
authentication error. No Aristotle project was created.
\end{itemize}

\section{Reproducibility}
Each file \texttt{proposals/P3.<ID>.lean} compiles with \texttt{lean} under toolchain
\texttt{leanprover/lean4:v4.28.0} and Mathlib commit
\texttt{8f9d9cff6bd728b17a24e163c9402775d9e6a365}. The probe table is regenerated by
\texttt{code/probe.sh proposals <scratch-dir>}. The rationale files
\texttt{proposals/P3.<ID>.md} list every hypothesis, every control and what each statement does
not settle. The two canon files used for comparison are copied from the canon repository at
commit \texttt{17a0702f}, and their hashes are recorded by the project.

\section{Changelog}
\begin{itemize}
\item 2026-11-05 (P3-T001): first version. Proposes P3.LoopMassBalance, P3.CascadedRecovery,
P3.RecoveryWorkFloor and P3.CyclicSymbiosis, with controls, probe results and rationale. No prior
P3 paper existed.
\end{itemize}

\bibliographystyle{plainnat}
\bibliography{refs}
\end{document}
